\uuid{gCNe}
\titre{Implication/équivalence}
\auteur{bodin ; hanani ; mzari}
\organisation{exo7}
\licence{CC BY-NC-SA 4.0}
\datecreate{2020-01-20}
\niveau{L1}
\module{Algèbre}
\chapitre{Logique, ensemble, raisonnement}
\sousChapitre{Absurde et contraposée}
\difficulte{3}
\theme{Raisonnement, implication, équivalence}
\exo7qcmid{lille-1.logique.029}

\contenu{
\begin{question}
\qclassification{100.03,100.04}
\qtags{motcle=implication/équivalence}
Quelles sont les assertions vraies ?
\begin{answers}
    \bad{La négation de "$P \implies Q$" est "non($Q$) ou $P$"}

    \good{La réciproque de "$P \implies Q$" est "$Q \implies P$"}

    \bad{La contraposée de "$P \implies Q$" est "non($P$) $\implies$ non($Q$)"}

    \bad{L'assertion "$P \implies Q$" est équivalente à "non($P$) ou non($Q$)"}
\end{answers}
\begin{explanations}
Il faut revenir à la définition de "$P \implies Q$" qui est "non($P$) ou $Q$".
\end{explanations}
\end{question}
}
